\section{Геометрия пути}\label{sec:geometry}

На протяжении статьи граф спуска \code{Step00} --- состояния (центры, дефекты, поглотители) и рёбра
(clean, boundary, peel, absorb) --- служил \emph{машиной}. Здесь мы читаем ту же конструкцию как
\emph{место} и задаём ей три вопроса, с которых началась геометрия: течёт ли время и необратимо ли
оно; какова форма пространства и можно ли её \emph{посчитать}, а не постулировать; пересекаются ли
прямые. Каждый ответ оказывается теоремой, читаемой с рёбер графа, а последний переворачивает самого
Евклида. Все результаты зелёные, кроме коды из двух деклараций; всю зелёную половину собирает капстоун
\code{geometry\_of\_euclids\_path}, не несущий пользовательских аксиом.

\subsection{Стрела времени}

Роль собственного времени играет \code{lexRank}, лексикографический ранг состояния --- не выбранный, а
вычисленный, строго падающий на каждом ребре.

\begin{theorem}[\code{timeArrow\_step}]\label{thm:timeArrow_step}
\green{} Каждый реальный шаг строго уменьшает ранг: из $\mathrm{RealStep}\ A\ M_0\ U\ V$ следует
$\code{lexRank}\,V < \code{lexRank}\,U$.
\end{theorem}

\emph{Идея доказательства.} Реэкспорт несущей теоремы конкретного графа: любой из четырёх
конструкторов шага --- clean, boundary, peel, absorb --- роняет ранг. Из этого локального факта
разворачивается вся кинематика: вдоль любого непустого пути ранг строго убывает
(\code{timeArrow\_path}), ни одно состояние не достижимо из самого себя (\code{no\_return}), а
бесконечного бега не существует вовсе (\code{no\_eternal\_run} --- прямое следствие корневого EPMI
\code{no\_infinite\_descent}, \cref{sec:engine}).

\begin{theorem}[\code{every\_journey\_halts}]\label{thm:every_journey_halts}
\green{} Любой маршрут останавливается не позже начального ранга: $k$ реальных шагов возможны только
при $k \le \code{lexRank}\,(\mathrm{run}\,0)$.
\end{theorem}

\emph{Идея доказательства.} Квалифицированный «второй закон», проведённый через
\code{Engine.turned\_engine\_halts}: строго убывающая на натуральных числах последовательность не
может сделать больше шагов, чем её начальное значение. Стрела времени --- теорема о конечном графе с
точной численной границей.

\subsection{Кривизна, вычисленная}

Метрики у нас нет --- есть только рёбра; но у ориентированного графа есть своя комбинаторная кривизна,
и она считается.

\begin{definition}[\code{curvature}]\label{def:curvature}
Кривизна вершины --- целое число
\begin{equation}\label{eq:curvature}
\kappa(v) = 1 - \operatorname{outdeg}(v),
\end{equation}
где $\operatorname{outdeg}$ --- число геодезических сегментов, выходящих из $v$ вниз.
\end{definition}

Это \emph{комбинаторная} кривизна --- дефицит или избыток исходящего потока геодезических, аналог
формулы Эйлера --- не риманова, не секционная, не Ricci--Ollivier: ни метрики, ни тензора, ни гладкой
структуры здесь нет, и слово «кривизна» законно лишь в смысле дискретной теории графов. Внутри этого
смысла всё честно: out-множество каждой вершины предъявлено как \code{Finset} и доказано совпадающим с
\code{RealStep} (\code{mem\_outTargets}), так что $\operatorname{outdeg}$ считает настоящие рёбра. Что
кривизна обязана смотреть вперёд --- тоже теорема.

\begin{theorem}[\code{inDegree\_infinite\_at\_origin}]\label{thm:inDegree_infinite_at_origin}
\green{} В исток \code{absorber 0} входит бесконечное семейство рёбер: каждый дефект
$\mathrm{defect}\ 0\ q\ \mathrm{minus}$ по всем $q$ втекает в могилу нуля.
\end{theorem}

\emph{Идея доказательства.} Absorb-ребро в \code{absorber 0} требует лишь $0 \le M_0$, а таких
дефектов бесконечно много. Значит, in-степень истока бесконечна, симметричная «полная степень» не
существует как \code{Finset}, и неориентированная кривизна невычислима в принципе; ориентированная
$\kappa = 1 - \operatorname{outdeg}$ --- единственная вычислимая версия.

При масштабе $(A, M_0) = (5, 4)$, каждое значение проверено ядром через \code{decide}, спектр есть
портрет мира. Поглотители несут $\kappa = +1$ (\code{curvature\_absorber}): полюса, где поток гаснет и
геодезические фокусируются. Дефекты с единственным выходом --- плоские коридоры $\kappa = 0$
(\code{curvature\_defect\_0\_2}, \code{curvature\_defect\_6\_5}). Дефект $(4,5,+)$, у которого есть и
peel, и absorb, ветвится при $\kappa = -1$ (\code{curvature\_defect\_4\_5}). Чистые центры --- всё
более гиперболические воронки: $\kappa = -3, -4, -8$ в центрах $2, 3, 7$ (\code{curvature\_center\_2},
\code{curvature\_center\_3}, \code{curvature\_center\_7}). Генеалогия гиперболична; дно сферично.
Дискретный Гаусс--Бонне стягивает картину.

\begin{theorem}[\code{gaussBonnet\_cone3}]\label{thm:gaussBonnet_cone3}
\green{} На forward-замкнутом световом конусе под центром $3$ полная кривизна равна эйлеровой
характеристике:
\begin{equation}\label{eq:gaussBonnet}
\chi(\mathrm{cone3}) = \textstyle\sum \kappa = -5.
\end{equation}
\end{theorem}

\emph{Идея доказательства.} Комбинаторное тождество $\sum\kappa = |W| - \sum\operatorname{outdeg}$
(\code{gaussBonnet\_sum}) на forward-замкнутом окне \code{cone3} (замкнутость проверена ядром,
\code{cone3\_forwardClosed}) даёт $-5$: два полюса и центр $0$ с $\kappa = +1$ каждый не перекрывают
ветвления воронок $-3$ и $-4$. Конус в целом гиперболичен, несмотря на сферическое дно. Нормировка
\cref{eq:curvature} --- выбор, и связь с настоящим Гаусс--Бонне --- аналогия, а не теорема
дифференциальной геометрии; но внутри неё характеристика посчитана точно, ядром, без единого допущения.

\subsection{Плоский мир --- вечный двигатель}

Могло ли пространство быть плоским всюду, $\kappa \equiv 0$? Не могло, и причина --- главный запрет
программы.

\begin{theorem}[\code{nonpositive\_curvature\_forces\_engine}]\label{thm:nonpositive_curvature_forces_engine}
\green{} Если бы кривизна была нигде не положительна ($\forall v,\ \kappa(v) \le 0$), из каждой
вершины шёл бы шаг, итерация выбора построила бы бесконечный бег --- запрещённый вечный двигатель.
Такого мира не существует.
\end{theorem}

\emph{Идея доказательства.} $\kappa(v) \le 0$ означает $\operatorname{outdeg}(v) \ge 1$: у каждой
вершины есть продолжение (\code{hasStep\_of\_curvature\_nonpos}). Если так всюду, выбор склеивает из
этих продолжений бесконечную геодезическую --- в точности вечный двигатель Евклида, убитый
\code{no\_eternal\_run}.

\begin{theorem}[\code{space\_positively\_curved\_somewhere}]\label{thm:space_positively_curved_somewhere}
\green{} Пространство пути Евклида где-то положительно искривлено: $\exists v,\ 0 < \kappa(v)$.
\end{theorem}

Кривизна вынуждена, а не декоративна: положительно искривлённая точка (полюс, дно, поглотитель) обязана
существовать, иначе спуск не имел бы дна и бежал бы вечно. Замкнутая геодезическая --- легальный
непустой цикл --- тоже была бы двигателем-свидетелем (\code{closedGeodesic\_is\_engineWitness}), уже
убитым ацикличностью \code{lexRank} (\code{no\_closedGeodesic}); плоскость и замкнутость --- два лица
одной запрещённой машины.

\subsection{Ирония постулатов}

Определим \emph{прямую} так, как её мыслил Евклид: максимальный путь спуска, продолжаемый, пока есть
куда шагнуть.

\begin{theorem}[\code{every\_line\_ends}]\label{thm:every_line_ends}
\green{} Каждая прямая упирается в терминал: из любого состояния за конечное число шагов доходишь до
вершины без выхода.
\end{theorem}

\emph{Идея доказательства.} Сильная индукция по собственному времени $\code{lexRank}\,v$: если $v$
терминально --- путь пуст; иначе первый шаг строго роняет ранг (стрела времени,
\cref{thm:timeArrow_step}), и хвост конечен по предположению.

\begin{theorem}[\code{no\_infinite\_line}]\label{thm:no_infinite_line}
\green{} Бесконечных прямых нет вовсе: никакую прямую нельзя продолжить неограниченно.
\end{theorem}

Вот ирония программы. «Прямую можно продолжить неограниченно» --- это \emph{второй} постулат
Евклида~\cite{euclid-elements}, скромнее и древнее знаменитого пятого. На пути Евклида падает именно он
(\code{euclid\_postulate\_two\_fails}): путь, названный его именем, нарушает его собственную аксиому ---
не пятую, о параллельных, а более глубокую, вторую, о самой продолжимости прямой. Параллельных тоже нет
(\code{no\_parallel\_lines}), но по вырожденной причине: параллельные суть две бесконечные прямые, а не
существует и одной. Здесь нужна скрупулёзность: это вырожденная, а не эллиптическая геометрия. Наивное
«все прямые встречаются» ложно.

\begin{theorem}[\code{bottom\_not\_unique}, \code{two\_disjoint\_lines}]\label{thm:two_disjoint_lines}
\green{} Дно не единственно --- \code{absorber 0} и \code{absorber 1} оба легальны и терминальны; и
существуют две полностью непересекающиеся конечные прямые (при $(5,4)$: маршрут
$\mathrm{center}\,7 \to \mathrm{defect}\,4\,5\,\mathrm{plus} \to \mathrm{absorber}\,4$ и маршрут
$\mathrm{center}\,3 \to \mathrm{defect}\,0\,2\,\mathrm{minus} \to \mathrm{absorber}\,0$), все девять
состояний которых попарно различны.
\end{theorem}

Так что «все линии сходятся в одну точку» в лоб --- неправда: есть две разные точки схода и пары
прямых, не встречающихся нигде. Читать падение второго постулата как «эллиптичность» законно лишь через
эту честную границу.

\subsection{Пересечение через могилу нуля}

Есть, тем не менее, честная, квалифицированная форма пересечения, проходящая через единственную особую
точку графа.

\begin{theorem}[\code{zeroPoint\_absorbs\_all\_divisors}]\label{thm:zeroPoint_absorbs_all_divisors}
\green{} Точка $0$ поглощает все делители: $\mathrm{sideValue}\ \mathrm{minus}\ 0 = 6\cdot 0 - 1 = 0$ в
$\mathbb{N}$, а на ноль делится всё; значит любое $q$ делит сторону нуля.
\end{theorem}

Это одновременно артефакт и маркер. Артефакт --- потому что усечённое вычитание в $\mathbb{N}$ даёт
$6\cdot 0 - 1 = 0$, тогда как в $\mathbb{Z}$ было бы $-1$; мы не утверждаем, что арифметика «знает»
это тождество вне $\mathbb{N}$-модели. Маркер --- потому что $0$ есть единственная точка, чьи стороны
поглощают все делители сразу, $\mathbb{N}$-след события $0 \to 1$ из первопричины --- начала стрелы
времени.

\begin{theorem}[\code{line\_to\_origin}]\label{thm:line_to_origin}
\green{} Из любого чистого центра $m \ge 1$ (при $2 \le A$) есть путь длины $2$ в могилу нуля:
boundary-шаг в дефект $(0,2,-)$ --- легальный, ибо $2 \mid 0$, --- затем absorb в \code{absorber 0}.
\end{theorem}

\begin{theorem}[\code{web\_above\_every\_horizon}]\label{thm:web_above_every_horizon}
\green{} Для $2 \le A$ и любого $N$ существует чистый центр $m > N$, чей путь ведёт в могилу нуля.
\end{theorem}

\emph{Идея доказательства.} Плотность не нужна: примориальный носитель
\code{Residuals.carrier\_nonempty\_above} предъявляет чистый центр выше любой границы (\cref{sec:carrier}),
а \code{clean\_of\_cleanZ} переводит его чистоту с $\mathbb{Z}$ на $\mathbb{N}$.

\begin{theorem}[\code{lines\_meet\_at\_origin}]\label{thm:lines_meet_at_origin}
\green{} Любые два чистых старта (при $2 \le A$) имеют общий терминал --- могилу нуля, \code{absorber 0}.
\end{theorem}

Не «единственная точка схода» и не «пересечение любых путей», а: у любой пары чистых линий есть общее
дно, к которому обе приходят. Все дороги вниз проходят через нуль. Пятый постулат опровергнут в честной
форме: прямые пересекаются --- на могиле нуля.

\subsection{Пересечение есть, но узнать его изнутри нельзя}

Остаётся эпистемический итог, и здесь единственные в разделе заражённые строки. Всё выше зелено; кода
несёт ровно две декларации, заражённые не через новую аксиому, а через уже принятую причинную границу
близнецов (\code{twinLowersInfinite\_from\_step00CausalClosure}), переведённую на язык вершин.

\begin{theorem}[\code{twin\_vertices\_beyond\_every\_horizon}]\label{thm:twin_vertices_beyond_every_horizon}
\yellow{} Над каждым горизонтом есть twin-вершина: для любого $N$ существует twin-центр $m > N$, обе
стороны которого $6m \pm 1$ просты. Условно на \axiom{}.
\end{theorem}

\emph{Идея доказательства.} Мост \code{twinCenter\_of\_twinLower} (сам зелёный, axiom-free) показывает
разбором вычета по модулю $6$, что младший член любой пары близнецов $p > 3$ есть $6m - 1$ для
twin-центра $m$; бесконечность близнец-младших --- ровно принятая граница. Никакого нового содержания
декрета: та же первопричина, что держит бесконечность близнецов, здесь лишь надета на вершины.

\begin{theorem}[\code{lines\_meet\_but\_unknowable\_from\_inside}]\label{thm:lines_meet_but_unknowable_from_inside}
\yellow{} Четырьмя фактами: (1) факт пересечения зелён; (2) twin-вершины есть над каждым горизонтом
(эта строка несёт таинт); (3) внутреннего основания факта пересечения нет; (4) будь оно --- это был бы
запрещённый двигатель Евклида.
\end{theorem}

\emph{Идея доказательства.} Факт пересечения \code{IntersectionFact} доказан зелено
(\code{intersectionFact\_green}: общий терминал есть могила нуля). Но вывод его как внутренней
первопричины \code{Step00}-мира --- \code{InternalisedIntersectionGround} --- по архитектуре есть
boundary-crossing self-proof, запрещённый двигатель. Его нет (\code{no\_internalisedIntersectionGround},
зелено), и знать его изнутри значило бы построить движок (\code{knowing\_meeting\_costs\_engine}). Более
того, это внутреннее основание тавтологично --- эквивалентно $P \wedge \neg P$
(\code{internalisedIntersectionGround\_iff}) --- так что всё содержание живёт в зелёном факте. Прямые
пересекаются, но узнать это изнутри нельзя: геометрический факт зелён; его внутреннее обоснование стоит
вечного двигателя.

\begin{remark}
Шесть границ нельзя раздувать. (i) $\kappa$ комбинаторна, не риманова. (ii) Гаусс--Бонне --- аналогия
при нормировке \cref{eq:curvature}, а не теорема дифференциальной геометрии. (iii) Могила нуля живёт на
$\mathbb{N}$-усечении: $6\cdot 0 - 1 = 0$ в $\mathbb{N}$, но $-1$ в $\mathbb{Z}$. (iv) Наивная
эллиптичность ложна; верна лишь квалифицированная \cref{thm:lines_meet_at_origin}. (v) Несущая теорема
отсутствия параллельных --- это падение \emph{второго} постулата (\cref{thm:no_infinite_line}), а не
пятого; геометрия вырожденная, не эллиптическая. (vi) Кода --- ровно две заражённые декларации, через
существующую границу близнецов, без новой аксиомы. Капстоун \code{geometry\_of\_euclids\_path} собирает
всю зелёную половину одной теоремой без пользовательских аксиом.
\end{remark}

Три древних вопроса получили ответы, и все три --- теоремы. Время течёт и необратимо; пространство
искривлено, и его кривизна посчитана ядром, с положительным дном и гиперболической генеалогией,
стянутой Гаусс--Бонне; прямые конечны --- падает второй постулат Евклида --- и все чистые линии
встречаются на могиле нуля, хотя узнать это изнутри нельзя. Получившаяся форма закручена внутрь:
гиперболическая генеалогия наверху, фокусирующее сферическое дно внизу, с $\chi(\mathrm{cone3}) = -5$,
тянущей к воронке. Это ставит путь Евклида рядом с теорией причинных множеств, где пространство-время
есть локально конечное частично упорядоченное множество --- наш вполне-упорядоченный ряд со строгой
стрелой и есть такой объект.
